\chapter*{序言}\label{chap:preface}
\markboth{序言}{序言}

\begin{quote}
最初的恐惧\\
是溺水, 所以\\
船最初的形状\\
是一只筏子; 等到\\
那事没发生, 要把它\\
弄得不再扁平可不容易...

\hfill ---Kay Ryan, \emph{We're Building the Ship as We Sail It}
\end{quote}

代数几何, 广义上是对多项式方程的几何研究, 如今已渗透到几乎整个现代数学. 二十世纪后半叶, 主要由于 Alexander Grothendieck 的工作, 这门学科经历了一场革命. 此后, 代数几何逐渐以一种近乎神话般强大, 却只有入门者才能看懂的机器而闻名. 许多优美的教材致力于逐步降低入门门槛, 帮助仅有一些交换代数和拓扑学基础的学生看清代数几何抽象概念背后的几何直观.

不过, 仍有一群学习者, 这门学科的大门尚未向他们完全开放: 对他们而言, 抽象代数的概念是一门新语言, 虽曾接触过, 却尚未完全内化. 有人可能坚持, 熟悉商环, 模和代数等概念是学习代数几何的先决条件. 我们则认为, 看到这些概念在几何情形中的表现, 可以巩固它们, 有时甚至可以由此引入它们. 这种视角除了让更多学生有机会进入代数几何, 还有一个优点: 代数学习者可以把注意力从任意环的广阔世界收回, 集中在多项式环这个更具体的范围内. 在这里, 许多概念具有更好的性质, 也更容易直观地理解它们的来由.

本书主要为这些学习者而写. 我们预期读者已经修过一门抽象代数入门课, 所以知道\textquotedbl{}环\textquotedbl{}和\textquotedbl{}理想\textquotedbl{}等词的含义. 不过, 读者也和我们一样是普通人: 例如, 见过商环的定义并不意味着就能毫无障碍地在两种看法之间切换, 即把商环视为陪集的集合, 或看成一种\textquotedbl{}置零\textquotedbl{}的操作. 我们力图在这些代数概念与所讲的几何发生关联时, 循序渐进, 自然而然地培养使用它们的能力.

另一方面, 我们也期待本书对代数基础更扎实, 希望学习代数几何基础的学生有所帮助. 对这样的学生而言, 书中介绍的许多代数思想或许并不新鲜, 但我们相信, 在几何情形中重新审视并巩固这些思想, 仍然值得.

\section*{预备知识}

如果你修过本科线性代数和抽象代数课程, 并且愿意深入思考, 慢慢阅读, 那么你已经具备学习本书的条件. 我们会不加定义地使用一些代数术语, 如环, 整环, 域和环同态. 环论中另外一些术语, 如多项式环和理想, 对我们展开代数几何尤其重要, 因此在 Chapter~\ref{chap:polynomial-rings} 中仔细定义. 我们建议所有读者, 尤其不确定自己代数基础是否足够的读者, 从细读该章开始. 在线性代数方面, 有时需要掌握一些矩阵计算, 包括行列式, 秩和零化度的基本性质; 我们也预期读者熟悉有限维向量空间, 包括张成空间, 线性无关, 基, 维数和线性映射等概念. 至于对偶空间和张量积等更深入的线性代数内容, 则在有需要时逐步展开.

值得注意的是, 我们没有把拓扑学基础作为先决条件. 尚未接触拓扑学的读者在本书 Part~\ref{part:affine-algebraic-geometry} (仿射代数几何) 不会处于劣势, 不过我们偶尔也会为学过一点拓扑学的读者指出相应联系. 在 Part~\ref{part:projective-algebraic-geometry} (射影代数几何) 中, 拓扑术语变得不可避免, 因此我们会在需要之处介绍相关背景. 即便这一部分也力求让没有拓扑学基础的学生能够读懂, 不过我们承认, 学过一门拓扑学课程肯定会使本书最后两章更容易消化.

\section*{如何使用本书}

本书沿着一条叙述主线展开, 因而适合按顺序阅读. 话虽如此, 有些内容可以自然地取出, 组成长度合适的课程. 例如, 一门专讲仿射代数几何的课程可以在一学期内讲完从 Chapter~\ref{chap:polynomial-rings} 到 Chapter~\ref{chap:products} 的大部分内容. 如果希望课程更快讲到射影簇, 也可以略过仿射代数几何中较深入的内容 (从 Chapter~\ref{chap:dimension} 到 Chapter~\ref{chap:products}, 讲的是维数, 光滑性和乘积), 直接跳到从 Chapter~\ref{chap:projective-varieties} 到 Chapter~\ref{chap:quasiprojective-varieties} 的内容. 对于代数背景较好的学生, 快速浏览部分纯代数内容之后, 也有可能在一学期内学完全书.

尽管本书有叙述性, 学生切不要以为可以像读小说一样阅读. 阅读时应手持铅笔, 常常停下来推敲细节, 演算例子. 目标是在细节理解和整体理解之间取得微妙的平衡: 熟悉细致计算, 同时不丢掉概念主线. 衡量自己做得如何的最佳办法, 是试着向同学清楚而简洁地解释所学内容.

说到这里, 我们想指出, 本书适合多种学习场景. 我们见过它成功地用于独立研习, 尤其是几名学生一起学习, 定期聚在一起讨论和解题的情形. 也可以用在翻转课堂里, 让学生至少通过课外阅读学习一部分内容. 即便完全独立自学, 我们也相信本书有用, 尤其对准备修读一门假定已熟悉簇的高阶代数几何课程的学生而言.

我们力求让本科生, 以及未受过代数几何训练但正在从事数学工作的读者, 都能读懂本书. 希望这本书能降低这门美丽学科的入门门槛, 向各个层次的学生和研究者打开大门, 让他们接触代数几何的严格基础.

\section*{体例选择}

为了尽可能提高可读性, 我们在体例和内容上作了不少选择. 本书符号几乎贯彻一致到了令人吃不消的程度. 例如, 我们用字母表末尾的字母表示变量, 用开头的字母表示变量的取值. 所以 $\mathbb{R}^{2}$ 中由 $y=x^{2}$ 定义的抛物线, 由满足 $b=a^{2}$ 的实数对 $(a,b)$ 组成. 这一选择源于我们的亲身体会: 对于不甚熟悉当前这类抽象表述的学生, 符号的一致性往往有帮助. 不过, 我们也知道, 最终应放下机械的一致性, 转而依靠灵活的直观; 我们试图让读者在这两端之间平稳过渡.

关于内容, 我们作了一个使本书区别于许多代数几何导论的重要决定: 在引入射影情形以前, 先比较完整地展开仿射簇的理论. 这个安排把代数与几何的相互作用放在中心位置, 因为它在仿射情形下最为显著; 同时, 它将最基础的代数进展与研究射影簇和拟射影簇所需的更多几何与拓扑概念分开. 不过, 如前所述, 如果特别想尽早熟悉射影簇, 也完全可以早一些跳到 Part~\ref{part:projective-algebraic-geometry}.

为写出一本篇幅合理的本科教材, 我们克制住了把自己最喜欢的所有题目都塞进书里的冲动, 选择在拟射影簇处结束, 恰好止于层, 抽象簇和概形论之前. 本书绝不是一部代数几何的综合参考书. 有些专家或许不同意我们舍去的内容, 不过我们已尽力写出连贯的叙述, 纳入我们认为属于代数簇严格基础的核心内容. 学完本书后, 学生将在拟射影簇的框架下打下坚实的基础, 从中汲取工具和直观, 继续学习专题或现代代数几何中更抽象的理论. 为引导后续学习, 书末的\hyperref[chap:coda]{后记}建议了若干最自然的下一步.
